\chapter{Conclusions}\label{chap:Conclusions}
This thesis has been an opportunity to investigate a new technique used in the search for New Physics.\\
The main section of this project was about our implementation of Anomaly Detection with a specific dataset and its transformations. The most interesting aspect is, for sure, the comparison between different types of dataset, which was presented to try to understand how a physical transformation can impact on performance of a computer learning. In fact, we started the thesis discussing about the "anti-$k_T$" dataset (the not transformed dataset) and it was crystal clear that, with an AUC score of 53\%, our model was acting as a random classifier. This failure was associated to rotational symmetries among data.\\
We decided to rotate data applying two Lorentz boosts in order to obtain $\eta'_{jet}=0$ and $\phi'_{jet}=0$ for each jet. This simple change resulted in a significant improvement in performance: we were able to get a 66.4\% in AUC score, which is not perfect, but at least our model is learning to discriminate between background and signal jets.\\
The third and last attempt was made transforming the dataset moving to a specific frame of reference. Performance is better than the rotated dataset, we got a 69\% AUC score, but, again, it is not perfect. \\
We can conclude that the transformed dataset is the one which gives us the best results, because the reduction of jet surface in the ($\eta , \phi $) plane makes the classification task easier for the model.\\
In the end, it is true that we were able to discriminate between background and signal jets but the quantitative parameter, the AUC score, underlined that this implementation is not perfect: 69\% of AUC in the best case scenario is not ideal if we want investigate small deviations from the Standard Model. This not brilliant result can be caused by not perfectly optimized hyperparameters or by the model itself. In fact, it is plausible that the architecture proposed does not fit very well for this task. Anyway, these results leave enough space for future development of the project. \\
We hope that one day, helped by Anomaly Detection, we will be able to find the church we are looking for, the New Physics that researchers have been investigating for years. 